\documentclass[11pt]{article}
\usepackage[margin=2.2cm]{geometry}
\usepackage{amsmath,amssymb,amsthm,mathtools}
\usepackage[most]{tcolorbox}
\usepackage[shortlabels]{enumitem}
\setlist{nosep}
\usepackage{fancyhdr,lastpage}
\usepackage[colorlinks=true,linkcolor=blue!50!black,urlcolor=blue!50!black]{hyperref}
\usepackage{xcolor}

\newcommand{\R}{\mathbb{R}}
\newcommand{\E}{\mathbb{E}}
\newcommand{\F}{\mathcal F}
\DeclareMathOperator{\Var}{Var}
\DeclareMathOperator{\Cov}{Cov}
\newcommand{\br}[1]{\langle #1\rangle}

\newtcolorbox[auto counter]{result}[2][]{%
  colback=blue!3,colframe=blue!45!black,
  title={Result~\thetcbcounter: #2},fonttitle=\bfseries,
  breakable,enhanced,boxrule=0.6pt,#1}
\newtcolorbox{sourcebox}{colback=gray!6,colframe=gray!55!black,
  boxrule=0.4pt,boxsep=2pt,left=6pt,right=6pt,top=3pt,bottom=3pt,
  fontupper=\small,breakable}
\newtcolorbox{trap}{colback=orange!5,colframe=orange!60!black,
  title=Common trap,fonttitle=\bfseries\small,fontupper=\small,breakable}

\newcommand{\used}[2]{\begin{sourcebox}\textbf{First appears:} #1 \\ \textbf{Reused in:} #2\end{sourcebox}}

\setlength{\headheight}{14pt}
\pagestyle{fancy}
\fancyhf{}
\fancyhead[L]{Mathematical Finance \& Stochastic Processes}
\fancyhead[R]{Recurring Results Toolbox --- Sheets 1--6}
\fancyfoot[C]{\thepage\ / \pageref*{LastPage}}
\renewcommand{\headrulewidth}{0.4pt}

\title{Recurring Results Toolbox\\[2pt]\large Mathematical Finance and Stochastic Processes in Continuous-Time\\ Collected from Problem Sheets 1--6 and their solutions}
\author{Kasper Svejgård Nielsen}
\date{\today}

\begin{document}
\maketitle
\thispagestyle{fancy}

\noindent This sheet collects results, techniques, and computational templates that
recur across Problem Sheets 1--6 (and their official solutions), with pointers to
every place they resurface. The goal is to stop re-deriving the same lemma from
scratch and instead recognise ``oh, this is \emph{Result~N} again.'' A typical
example: Sheet 4, Problem 4.3(h) derives $\Var(\int_s^{s+t}B_r\,dr)=st^2+t^3/3$ to
show the integrated Brownian motion does \emph{not} have stationary increments --
the underlying technique (Fubini + $\E[B_rB_s]=r\wedge s$, then Theorem~4.8 to rule
out the martingale property) is exactly the one needed for the covariance
computations in Problem~6.3 and for several quiz-style questions. That technique is
catalogued below as Result~16 (the ``not a martingale'' test) and
Result~17 (the time-integral moment formulas).

\tableofcontents

%======================================================================
\section{Stopping times: construction and recognition}
%======================================================================

\begin{result}{Hitting times of continuous adapted processes are stopping times}
\label{res:1}
If $X=(X_t,t\ge0)$ is continuous and adapted and $A\subseteq\R$ is closed, then
$T=\inf\{t\ge0\mid X_t\in A\}$ is a stopping time (Proposition 3.9(ii)). In
particular, hitting times of a fixed level, or of a moving but continuous
deterministic boundary $b(t)$, are stopping times: just apply the result to
$X_t - b(t)$.
\end{result}
\used{Problem 2.1(d) ($T\wedge t$ argument via Prop.\ 3.6); Problem 3.4(a)
($T=\inf\{t\mid B_t=t-1\}$, via $X_t=B_t-t$).}{Quiz practice problems on
stopping-time identification (Quizzes A/B, Q1).}

\begin{result}{Last-exit / ``time of the extremum'' is \emph{not} a stopping time}
\label{res:2}
Random times defined as a $\sup$ of a future-dependent set --- last passage times
$\rho=\sup\{t\ge0\mid X_t\ge a\}$, or the time $\theta$ at which a running maximum
is attained on a fixed horizon --- are \emph{not} stopping times in general: the
event $\{T\le t\}$ requires knowing the path strictly after time $t$.
\end{result}
\used{Problem 2.2(b) (last passage time above level $a$).}{Problem 3.4(b) (time of
the maximum of $B$ on $[0,1]$); same reasoning re-used verbatim in solution 3.4(b).}

\begin{result}{Building new stopping times from old ones}
\label{res:3}
\begin{itemize}
\item For $A\in\F_s$, the time $T^A(\omega)=s\cdot\mathbf 1_A(\omega)+\infty\cdot
\mathbf 1_{A^c}(\omega)$ is a stopping time.
\item More generally, for $A\in\F_s$ and $s<t$, $T(\omega)=s$ on $A$ and $t$ on
$A^c$ is a stopping time (special case: $T=T^A\wedge t$).
\item Minima, maxima, and sums of stopping times are stopping times; the minimum
of a stopping time and a deterministic constant is a stopping time.
\end{itemize}
\end{result}
\used{Problem 2.2(a) ($T^A$); page-45 property list.}{Problem 5.2(a) (re-derives
$T^A\wedge t$ is a stopping time two ways); Problem 5.3(b) (localizing sequences
$\tilde T_n$ built from the same $T^A$-template).}

%======================================================================
\section{The exponential-martingale engine}
%======================================================================

\begin{result}{MGF-as-ODE trick}
\label{res:4}
If $M(\lambda)=\E[e^{\lambda X}]$ satisfies $M'(\lambda)=\lambda\sigma^2 M(\lambda)$
with $M(0)=1$, define $H(\lambda)=e^{-\lambda^2\sigma^2/2}M(\lambda)$. Then
$H'\equiv0$, so $H\equiv1$ and $M(\lambda)=e^{\lambda^2\sigma^2/2}$: this
\emph{is} the defining computation behind every ``exponential process is a
martingale'' argument in the course.
\end{result}
\used{Problem 1.1(d)--(e) (MGF of $N(0,1)$).}{The identical ODE argument reappears,
dressed up as a conditional-expectation computation, in Problem 3.1(b), Problem
4.1(i), and Problem 6.2(b)(i) -- see Result~7.}

\begin{result}{Independence from a factorised joint MGF / characteristic function
(Problem 1.2, Lemmas 1--2)}
\label{res:5}
If $\E[e^{\lambda X}\mid\mathcal G]=e^{\lambda^2\sigma^2/2}$ for all $\lambda$ (or
the characteristic-function analogue), then $X\sim N(0,\sigma^2)$ \emph{and} $X$
is independent of $\mathcal G$. Proof template: tower property $+$ pull-out
property turns the conditional MGF into a product of marginal MGFs/characteristic
functions, which characterises independence.
\end{result}
\used{Problem 1.2 (full proof, both the MGF and characteristic-function
versions).}{Problem 6.2(c)(i): applied directly to $M_t-M_s$ conditional on
$\F_s$ to show a local martingale with deterministic quadratic variation has
independent Gaussian increments -- this is exactly Lévy's characterisation
specialised to deterministic $\br{M,M}$.}

\begin{result}{Monotone-class extension of independence (Problem 1.3)}
\label{res:6}
$(X_i,i\in I)$ is independent of $\mathcal G$ iff every \emph{finite} subcollection
$(X_{i_1},\dots,X_{i_p})$ is independent of $\mathcal G$. This lets you check
independence on finite-dimensional marginals and automatically lift it to an
entire (possibly infinite) family / $\sigma$-algebra.
\end{result}
\used{Problem 1.3 (statement and proof via the monotone class theorem).}{Problem
2.3(c)(ii): used twice, first to go from finite collections of increments to
$\sigma(B_r\mid 0\le r\le s)$, then to lift to the whole process $(X_t,t\ge0)$.}

\begin{result}{\textbf{The} exponential-martingale template}
\label{res:7}
Let $M$ be a continuous local martingale with $M_0=0$ and deterministic quadratic
variation $\br{M,M}_t=g(t)$. Then for every $\theta\in\R$,
\[
Y_t=\exp\!\Big(\theta M_t-\tfrac12\theta^2 g(t)\Big)
\]
is a martingale. This single statement specialises to \emph{every} ``find
$\theta$ such that $\exp(\cdots)$ is a martingale'' question in the course.
\end{result}
\begin{itemize}[leftmargin=*]
  \item $M=B_t$, $g(t)=t$: Problem 3.1(b) ($Y_t=e^{\theta B_t-\theta^2t/2}$), and
  Problem 3.1(c) shows the $\mu\ne0$ drift version is a sub-/supermartingale.
  \item $M=\boldsymbol\lambda\cdot\boldsymbol B_t$ in $\R^d$ with
  $\boldsymbol\lambda\cdot\boldsymbol\lambda=1$: Problem 3.2(c).
  \item $M$ a general continuous Gaussian martingale, $g(t)=\Gamma(t,t)$: Problem
  4.1(i), used immediately in the quiz question 4.2(b) ($\theta=\pm2$ for
  $g(t)=2t$ against a target coefficient of $4t$).
  \item $M$ a continuous local martingale with \emph{any} deterministic,
  continuous, increasing quadratic variation $g(t)=\br{M,M}_t$ (not assumed
  Gaussian or Brownian a priori): Problem 6.2(b)(i)--(ii), which then
  \emph{bootstraps Result~7 to prove $M$ is Gaussian in the first place}
  (Problem 6.2(c)--(d)) -- the martingale property comes first, Gaussianity is a
  consequence, not an assumption.
\end{itemize}
\used{Problem 3.1(b).}{Problems 3.1(c)(d), 3.2(c), 4.1(i), 4.2(b)(quiz), 6.2(b),
and implicitly Problem 5.5(a) (showing the \emph{compensator cannot be dropped}:
$e^{\lambda B_t}$ alone is a martingale only for $\lambda=0$).}

\begin{trap}
Problem 5.5(a) is the ``negative'' companion to Result~7: drop the
$-\theta^2g(t)/2$ compensator and the process is a martingale \emph{only} for
$\theta=0$. Always check whether a given exponential process already includes
the compensator before answering a ``for which $\theta$'' question.
\end{trap}

%======================================================================
\section{Computing moments of Brownian / Gaussian-martingale functionals}
%======================================================================

\begin{result}{``Pull out the past, the future has mean zero'' rule}
\label{res:8}
For $0\le s<t$ and any $\F_s$-measurable $Z$ with $\E|Z|<\infty$,
\[
\E[Z\cdot(B_t-B_s)]=\E[Z]\cdot\E[B_t-B_s]=0,
\qquad
\E[Z\cdot g(B_t-B_s)]=\E[Z]\cdot\E[g(B_t-B_s)]
\]
whenever $B_t-B_s$ is independent of $\F_s$ (true for Brownian motion, and for any
process with independent increments, e.g.\ Problem 4.1's Gaussian martingales).
This is the single most-used computational step in the course.
\end{result}
\used{Problem 3.1(a) (expanding $B_t^2-t$ via $B_t=B_s+(B_t-B_s)$).}{Problem
3.4(c)--(e) (all four moment computations); Problem 4.1(e); quiz 4.2(c); Problem
5.5(b)(i); and essentially every ``compute $\E[\cdots]$'' quiz question on
Sheets 3--6.}

\begin{result}{Odd moments of centred (jointly) Gaussian variables vanish}
\label{res:9}
If $X\sim N(0,\sigma^2)$ then $\E[X^{2n+1}]=0$ for all $n\ge0$. More generally,
any odd-order moment of a mean-zero jointly Gaussian vector
$(X_1,\dots,X_k)$ -- e.g.\ $\E[X_1X_2X_3]$ or $\E[X_1^2X_2]$ -- is zero, by the
symmetry $(X_1,\dots,X_k)\overset{d}=(-X_1,\dots,-X_k)$.
\end{result}
\used{Problem 1.1(c) (derived by induction via Stein's identity).}{Problem
3.4(d)(f) ($\E[B_s^2B_t]=0$, $\E[B_sB_tB_u]=0$ -- solution sheet gives \emph{two}
independent derivations, one via Result~8 and one via the tower property, as a
cross-check); quiz questions throughout Quizzes A/B.}

\begin{result}{Covariance shortcut for Brownian motion and Gaussian martingales}
\label{res:10}
$\Cov(B_s,B_t)=s\wedge t$. More generally, if $M$ is a continuous Gaussian
martingale with $M_0=0$ and $\Cov(M_s,M_u-M_t)=0$ for $s\le t<u$ (Problem
4.1(e)), then for $s\le t$,
\[
\Cov(M_s,M_t)=\Gamma(s,s)=\Gamma(s\wedge t,\,s\wedge t).
\]
\end{result}
\used{Problem 4.1(e) (proved via tower property + martingale property).}{Problem
3.3(b)(i) (direct analogue for the correlated 2D Brownian motion,
$\Cov(W^{(1)}_s,B^{(2)}_t)=\rho(s\wedge t)$); quiz 4.2 / 6.3(g)
($\Cov(B_t,Y_t)$ computed via the Itô-isometry version, Result~12).}

%======================================================================
\section{Quadratic variation, Itô integrals, and the product rule}
%======================================================================

\begin{result}{Bracket bilinearity for stochastic integrals}
\label{res:11}
For Itô integrals driven by the same (or independent) Brownian motions,
\[
\Big\langle\!\int H\,dB^{(1)},\int K\,dB^{(2)}\Big\rangle_t
=\int_0^t H_sK_s\,d\br{B^{(1)},B^{(2)}}_s,
\]
and $\br{B^{(1)},B^{(2)}}_t=0$ whenever $B^{(1)}\perp B^{(2)}$ are independent
Brownian motions (Proposition 4.16). When $B^{(1)}=B^{(2)}=B$, this collapses to
$\br{\int H\,dB,\int K\,dB}_t=\int_0^t H_sK_s\,ds$.
\end{result}
\used{Problem 5.4(b)--(c) (orthogonality of integrals against independent
Brownian motions).}{Problem 6.1 (every sub-part: computing $\br{W^{(i)},W^{(j)}}_t$
to invoke Lévy's characterisation); Problem 6.3(a)(b) ($\br{Y,Y}_t=\int s^2ds$,
$\br{Y,Z}_t=\int sY_s\,ds$); Problem 6.4(a); quiz questions on $\br{Y,Y}_t$.}

\begin{result}{Deterministic / bounded integrand $\Rightarrow$ genuine
martingale, via Itô isometry}
\label{res:12}
If $H$ is adapted and $\E\int_0^t H_s^2\,ds<\infty$, then $X_t=\int_0^t H_s\,dB_s$
is a (true, not just local) square-integrable martingale with
\[
\E[X_t]=0,\qquad \Var(X_t)=\E[X_t^2]=\E\Big[\int_0^t H_s^2\,ds\Big].
\]
\end{result}
\used{Problem 5.5(c)(i), 5.6(c)(d) (variance of $\int_0^tB_s\,dB_s$).}{Problem
5.4(d) (used in $\mathbb H^2$ form to get $\E[X_tY_t]=0$); Problem 6.3(c)(d);
Problem 6.4(a)(b); every ``is this an Itô integral a martingale'' quiz question.}

\begin{trap}
The martingale property in Result~12 needs $H$ to not depend on the \emph{upper
limit} $t$ itself. $X_t=\int_0^te^{-(t-s)}\,dB_s=e^{-t}\int_0^te^sdB_s$ looks like
a stochastic integral but is \emph{not} one in the sense of Result~12: writing
$Z_t=\int_0^te^sdB_s$ (a genuine martingale) and $X_t=e^{-t}Z_t$, the product rule
gives $dX_t=-X_t\,dt+dB_t$ -- a nonzero drift appears, so $X$ is \emph{not} a
martingale (it is the stationary Ornstein--Uhlenbeck process). Contrast with
Problem 5.5(c)(ii), which is exactly this trap.
\end{trap}

\begin{result}{$M^2-\br{M,M}$ is a martingale (and its converse: identifying
$\br{M,M}$)}
\label{res:13}
For a continuous (local) martingale $M$ with $M_0=0$, $M_t^2-\br{M,M}_t$ is a
martingale. Conversely, since $\br{M,M}$ is the \emph{unique} continuous
increasing process making $M^2-\br{M,M}$ a martingale, if you can show
$M_t^2-g(t)$ is a martingale for some deterministic, continuous, increasing
$g$, then necessarily $\br{M,M}_t=g(t)$.
\end{result}
\used{Problem 3.1(a) ($B_t^2-t$, base case).}{Problem 3.2(b) (multivariate sum of
base cases); Problem 4.1(g)--(h) (general Gaussian martingale, $g=\Gamma(t,t)$,
used to \emph{identify} the quadratic variation rather than just verify it);
quiz 4.2(a); Problem 6.2(a).}

\begin{result}{Product-minus-bracket is a local martingale; orthogonality}
\label{res:14}
For continuous local martingales $M,\tilde M$, the process $M\tilde M-\br{M,\tilde
M}$ is a local martingale. Consequently $M\perp\tilde M$ (i.e.\ $\br{M,\tilde
M}\equiv0$) iff $M\tilde M$ is itself a local martingale.
\end{result}
\used{Problem 5.4(a) (statement, via Proposition 4.15(i)).}{Problem 3.3(c) (the
\emph{same} computation carried out by hand, one week earlier, for
$M_t=W^{(1)}_tB^{(2)}_t-\rho t$ -- a good worked example of Result~14 before the
general machinery is available); Problem 6.3(f) ($\E[Y_tZ_t]=\E[\br{Y,Z}_t]$).}

%======================================================================
\section{Gaussian-process closure, finite variation, and time integrals}
%======================================================================

\begin{result}{Gaussianity is closed under linear maps and $L^2$/pointwise limits
(Proposition 1.1)}
\label{res:15}
A finite linear combination of jointly Gaussian variables is Gaussian; and if
$X_n\to X$ in $L^2$ (or a.s., or in distribution) with each $X_n$ Gaussian, then
$X$ is Gaussian. This is the standard route to proving a \emph{limit} process is
Gaussian without computing its density directly.
\end{result}
\used{Problem 2.6(a) (Brownian bridge is Gaussian, as a linear map of
$(B_{t_1},\dots,B_{t_p},B_1)$).}{Problem 4.3(b)--(c) (integrated Brownian motion:
Riemann sums are Gaussian, limit is Gaussian); Problem 6.3/6.4(c)--(e) (elementary
functions $\to$ Gaussian processes, approximating $\int f\,dB$).}

\begin{result}{A continuous martingale of finite variation is constant
(Theorem~4.8) --- the standard ``NOT a martingale'' test}
\label{res:16}
If $X$ is a continuous martingale and has (a.s.) finite variation, then $X$ is
a.s.\ constant. Contrapositive: if $X_0$ is deterministic, $X$ is continuous and
of finite variation, and $\Var(X_t)\ne0$ for some $t$, then $X$ is \emph{not} a
martingale.
\end{result}
\used{Problem 4.3(e): the integrated Brownian motion $\int_0^tB_s\,ds$ is
continuous, has finite variation (Problem 4.3(c)), and $\Var(\int_0^tB_s\,ds)=
t^3/3\ne0$ -- hence it cannot be a martingale, even though it is centred
(Problem 4.3(d)).}{This is the standard argument whenever a process looks
``martingale-like'' (centred, continuous) but is built from a time-integral
rather than a stochastic integral. Watch for it on any question of the form
``is $\int_0^t X_s\,ds$ a martingale?''}

\begin{result}{Second moments of time-integrals of Brownian motion (Fubini +
$\E[B_rB_s]=r\wedge s$)}
\label{res:17}
For $0\le s<t$, writing $I_t=\int_0^tB_r\,dr$:
\[
\E[I_t]=0,\qquad
\Var(I_t)=\frac{t^3}{3},\qquad
\Cov(I_s,I_t)=s^2\Big(\frac t2-\frac s6\Big).
\]
Consequently the increments are \emph{not} stationary:
$\Var(I_{s+t}-I_s)=st^2+\dfrac{t^3}{3}\ne\dfrac{t^3}{3}=\Var(I_t)$,
and \emph{not} independent, e.g.\ $\Cov(I_1,I_2-I_1)=\tfrac12\ne0$.
\end{result}
\used{Problem 4.3(d), (f)(ii), (g), (h) (full derivation).}{Problem 4.4(b)(f)(g)
(simulation exercise checks these exact numbers numerically); the underlying
technique (Fubini + the covariance kernel $r\wedge s$ of $B$) is the direct
template for Problem 6.3(c)--(g)'s moment computations for $Y_t=\int_0^tsdB_s$
and $Z_t=\int_0^tY_s\,dB_s$, and for quiz questions computing
$\Cov(B_t,Y_t)=\int_0^ts\,ds=t^2/2$.}

\begin{result}{Toolbox of Brownian-motion-preserving transformations}
\label{res:18}
Each of the following is again a standard Brownian motion whenever $B$ is:
\begin{itemize}
\item $-B_t$ \ (symmetry);
\item $\lambda B_{t/\lambda^2}$ for $\lambda>0$ \ (Brownian scaling);
\item $B_{s+t}-B_s$ for fixed $s\ge0$, and this new process is independent of
$\F^B_s$ \ (time-homogeneity / strong Markov flavour);
\item $W_t=tB_{1/t}$ for $t>0$, $W_0=0$ \ (time inversion);
\item $\boldsymbol\lambda\cdot\boldsymbol B_t$ for $\boldsymbol\lambda\in\R^d$
with $\boldsymbol\lambda\cdot\boldsymbol\lambda=1$ \ (unit-vector projection of
multivariate BM).
\end{itemize}
\end{result}
\used{Problem 2.3 (first four transformations, each proved from the
definition).}{Problem 2.4(b) (time inversion, used together with Result~18's
SLLN consequence $B_t/t\to0$ a.s.\ from Problem 2.4(a)); Problem 3.2(a) (unit
projection); Problem 3.3(a) (built directly from the projection result with
$\boldsymbol\lambda=(\sqrt{1-\rho^2},\rho)$); underlies every ``is this
combination of stochastic integrals a Brownian motion'' question in Problem 6.1
(via Lévy's characterisation instead, but motivically the same idea).}

%======================================================================
\section{Supporting inequalities and convexity facts}
%======================================================================

\begin{result}{Doob's $L^2$ maximal inequality}
\label{res:19}
For a continuous martingale $M$, $\E\big[\sup_{t\in[0,T]}M_t^2\big]\le4\,\E[M_T^2]$.
\end{result}
\used{Problem 2.4(a)(iv) (bounds $\E[S_n]<\infty$ for the sup-of-increments
sequence, needed for the SLLN step).}{Problem 4.1(d) (dominated convergence
argument for continuity of $t\mapsto\Gamma(t,t)$).}

\begin{result}{Convex functions of a martingale are submartingales (Proposition
3.12)}
\label{res:21}
If $M$ is a martingale and $h$ is convex (and, when needed, increasing in the
relevant direction), then $h(M_t)$ is a submartingale.
\end{result}
\used{Problem 3.1(e): $(e^{B_t}-K)^+$ is a submartingale, since $x\mapsto
(e^x-K)^+$ is convex (composition of convex $e^x-K$ with convex increasing
$y\mapsto y^+$).}{Problem 4.1(d): $M_t^2$ is a submartingale since $x\mapsto x^2$
is convex -- used to show $t\mapsto\Gamma(t,t)=\E[M_t^2]$ is increasing.}

%======================================================================
\section{Multivariate It\^o, the product rule, the optional stopping theorem,
and the reflection principle}
%======================================================================

\noindent The results below are not stated explicitly on Problem Sheets 1--6,
but they are standard consequences of machinery that \emph{is} developed there
(It\^o's formula, the martingale property, L\'evy's characterisation). Quiz D
is the first practice quiz to require them directly, so they are collected here
rather than being re-derived from scratch each time.

\begin{result}{It\^o's formula for a function of several independent Brownian
motions}
\label{res:22}
Let $B^{(1)},\dots,B^{(d)}$ be independent Brownian motions and let
$f(x_1,\dots,x_d,t)$ be $C^2$ in $x$ and $C^1$ in $t$. Then
\[
df\big(B^{(1)}_t,\dots,B^{(d)}_t,t\big)=\sum_{i=1}^d f_{x_i}\,dB^{(i)}_t
+\Big(f_t+\tfrac12\sum_{i=1}^d f_{x_ix_i}\Big)dt,
\]
i.e.\ exactly the single-variable It\^o formula applied coordinate-by-coordinate,
with \emph{no} cross terms $f_{x_ix_j}$ ($i\ne j$) because
$\langle B^{(i)},B^{(j)}\rangle_t=0$ for independent Brownian motions (Toolbox
Result 11). Contrast with Problem 3.3(c) / Toolbox Result 14, where the two
driving Brownian motions are \emph{correlated} and a cross term
$\rho\,\partial^2_{xy}f\,dt$ survives.
\end{result}
\used{Problem 3.1(a) (single-variable case, $f(x,t)=x^2-t$).}{Quiz D, Q2
($f(x,y)=x^2+y^2$, identifying the compensator $c=2$).}

\begin{result}{It\^o's product rule, in stochastic-differential form}
\label{res:23}
For continuous (local) martingales (or more general It\^o processes) $X,Y$,
\[
d(X_tY_t)=X_t\,dY_t+Y_t\,dX_t+d\langle X,Y\rangle_t.
\]
This is the differential form of ``product minus bracket is a (local)
martingale'' (Toolbox Result 14): integrating $d(XY)$ term by term reproduces
$X_tY_t-\langle X,Y\rangle_t$ as the martingale part. It is the fastest way to
find the SDE satisfied by a product of two known It\^o processes, without
re-deriving the bracket from first principles each time.
\end{result}
\used{Problem 3.3(c) ($d(W^{(1)}_tB^{(2)}_t)$, derived by hand before the general
rule is available).}{Problem 5.4(a) (general statement, via Proposition 4.15(i));
Quiz D, Q6 (drift of $Z_t=B_tX_t$ for $X_t=\int_0^ts\,dB_s$, using
$\langle B,X\rangle_t=t^2/2$ from bracket bilinearity, Toolbox Result 11).}

\begin{result}{Optional stopping theorem (bounded / uniformly integrable case)}
\label{res:24}
If $M$ is a continuous martingale and $T$ is a stopping time such that the
stopped process $(M_{t\wedge T})$ is uniformly integrable (e.g.\ $T$ is bounded,
or $T$ is a.s.\ finite with $\sup_t\E[M_{t\wedge T}^2]<\infty$), then
$\E[M_T]=\E[M_0]$. Applying this to \emph{both} $M_t=B_t$ and $M_t=B_t^2-t$
(Toolbox Result 13) at the same exit time $T$ is the standard ``gambler's ruin
for Brownian motion'' template: the first identity pins down the exit
probabilities, the second then gives $\E[T]$.
\end{result}
\used{Problem 2.1(d) (bounded stopping time $T\wedge n$, used to control a
different quantity).}{Quiz D, Q4--Q5: $T=\inf\{t\mid B_t\in\{3,-2\}\}$;
$\E[B_T]=0$ gives $P(B_T=3)=2/5$ (Q4), then $\E[B_T^2-T]=0$ gives $\E[T]=6$
(Q5).}

\begin{trap}
Optional stopping requires \emph{some} integrability or boundedness condition on
$T$ -- it is not automatic just because $T<\infty$ a.s. (The classical
counterexample: for $T=\inf\{t\mid B_t=1\}$, $\E[B_T]=1\ne0=\E[B_0]$, because
$T$ fails the uniform-integrability condition.) Quiz D's Q4--Q5 sidestep this by
stating explicitly that the theorem may be assumed to apply; always check this
hypothesis before invoking the theorem elsewhere.
\end{trap}

\begin{result}{Reflection principle}
\label{res:25}
For $a\ge0$ and $t>0$,
\[
P\Big(\max_{s\in[0,t]}B_s\ge a\Big)=2\,P(B_t\ge a)=2\big(1-\Phi(a/\sqrt t)\big),
\]
equivalently $\max_{s\in[0,t]}B_s\overset d=|B_t|$. This converts a question
about the running maximum of $B$ into a single Gaussian tail probability, and
underlies the hitting-time distribution of $T_a=\inf\{t\mid B_t\ge a\}$ via
$\{T_a\le t\}=\{\max_{s\in[0,t]}B_s\ge a\}$.
\end{result}
\used{Problem 3.5(c) (simulation: histogram of $\max_{[0,1]}B_s$ compared to the
density of $|B_1|$ -- the simulation-based version of this identity).}{Quiz D,
Q3 ($P(\max_{[0,4]}B_s\ge2)=2(1-\Phi(1))$).}

\begin{result}{Conditional distribution of Brownian increments (Markov
property)}
\label{res:26}
For $0\le s<t$, $B_t\mid\mathcal F_s\sim N(B_s,t-s)$, because
$B_t-B_s\sim N(0,t-s)$ is independent of $\mathcal F_s$. Consequently, for any
Borel set $A$, $P(B_t\in A\mid\mathcal F_s)=P(B_s+Z\in A)$ with $Z\sim
N(0,t-s)$ -- in particular, since $B_t-B_s$ is symmetric about $0$,
$P(B_t>B_s\mid\mathcal F_s)=1/2$ regardless of $B_s$.
\end{result}
\used{Problem 2.1 (filtration/Markov background).}{Quiz D, Q10 (the symmetry
shortcut $P(B_t>B_s\mid\mathcal F_s)=1/2$, a common trap for students who
instead reach for $\Phi(B_s/\sqrt{t-s})$).}

\vspace{8pt}
\begin{center}\rule{0.5\textwidth}{0.4pt}\\[2pt]\textbf{END}\end{center}

\end{document}
